%=============================================================================
% DATASET -- a contribution, so described as method (Madam's rule).
% Budget: about half a column. Double-blind: no "our earlier work".
%=============================================================================
\section{Dataset Construction}
\label{sec:data}

\textbf{Auditing the benchmark.} The array used by the existing
benchmark~\cite{weng2024graph} has no date axis. Dating each row by matching its 25
district counts against the published Weekly Epidemiological Reports (WER) shows that
rows 0--47 are from 2023 and rows 51 onward run 2013--2022, so 45 rows from 2023 sit in
the training split of every fold that tests on 2018--2022. The weather channels were
not reordered with the cases: the array's temperature matches ERA5 at $r=0.95$ on its
own row order but $r=0.45$ on the dates the case rows carry, and its ``lag-12''
rainfall matches ERA5 rain twelve weeks \emph{later} ($r=0.68$, against $0.19$ twelve
weeks earlier). Its largest spike, a 19-fold jump in one week, is a dragged
spreadsheet formula in one printed table (15 of 24 cells equal the sum of the two
before). Two of these findings leak future information, so scores on this array are
optimistic.

\textbf{The corrected dataset.} We rebuild the series from original files, each checked
against the SHA-256 recorded at first retrieval: the dengue rows of every WER, ERA5
daily weather, and census population. The result covers $N=25$ districts over 559 weeks
(2013-W26 to 2024-W10), with weekly cases, temperature, precipitation, humidity, soil
moisture and population. Seven weeks with no report stay missing and are never
interpolated. A forecast for week $t$ reads cases up to $t-1$ and climate up to $t-2$
(ERA5 is released about five days late), and a perturbation test asserts that no later
value can change a forecast. The district graph is the benchmark's adjacency list with self-loops (141 directed
entries); the prospective test uses a corrected version with two non-border entries
removed (139).

\textbf{Extension to 2026.} For a prospective test we parse the 127 WERs from 2024-W11 to
2026-W32 under quality rules committed before any count was read: district totals must
equal the national total, and a dragged-formula pattern or a reprinted table voids a
week, which is replaced by the difference of the cumulative rows or left missing. This
gives 109 observed, 8 corrected and 10 missing weeks. A first scoring run exposed two
reconstruction defects (the first report of a year carries the previous year's total,
and two reports reprint the previous week's table). The rules were amended on report
semantics alone, with no model setting changed, so we call this an \emph{amended}
prospective evaluation.
